\begin{frame}{경로 = 설명: 결정 트리의 강점}
  \begin{itemize}
    \item 새 샘플 들어오면: (1) 루트에서 질문을 평가 $\to$ (2) 예/아니오에
      따라 좌·우 자식으로 이동 $\to$ (3) 리프면 그 값을 내놓고 끝
    \item 이 \textbf{경로 자체가} ``왜 이렇게 예측했는가''의 설명이다
    \item 대출 심사 예: ``소득 $<$ 3000만원 $\to$ 연령 $>$ 45세 $\to$
      \textbf{거절}''이라는 경로를 \textbf{그대로 읽어}주면 됨
  \end{itemize}
  \vskip0.8em
  \begin{block}{\textbf{설명가능성}(interpretability)}
    루트에서 리프까지의 질문 경로를 사람이 읽을 수 있다. 이 장(7.3)에서는
    \textbf{수백 개의 트리}가 합쳐지면 이 ``경로 = 설명'' 특성을 잃는
    대가를 다시 논한다(SHAP).
  \end{block}
\end{frame}
